\chapter{Cost--Benefit Analysis}\label{ch:b3:cost-benefit-analysis}

A government deciding whether to build a bridge, regulate a chemical or fund a
programme needs a common unit for what it costs and what it gives. \omterm{def:b3:cost-benefit-analysis:cba}{Cost--benefit
analysis} supplies one: money, measured by what people would be willing to pay.
This chapter states the method, measures the two numbers on which most analyses
turn --- the discount rate and the \omterm{def:b3:cost-benefit-analysis:vsl}{value of a statistical life} --- and sets the
forecasts of large projects against their outcomes.

\begin{openingchart}{Average cost overrun of 258 transport projects, actual against
estimated construction cost at fixed prices, by type of project and region
(percent; no data for fixed links and roads outside Europe and North
America).}{Flyvbjerg, Holm and Buhl (2002), Tables 1 and 2}
\begin{tikzpicture}
\begin{axis}[omrecord, ybar, bar width=7pt, width=0.66\linewidth, height=0.4\linewidth,
  xmin=-0.6, xmax=3.6, ymin=0, ymax=70, xtick={0,...,3},
  xticklabels={Rail, Fixed links, Roads, All projects}, ylabel={Overrun (\%)},
  grid=none, ymajorgrids, area legend, legend style={at={(0.5,-0.18)},
  anchor=north, legend columns=4, font=\footnotesize}]
  \addplot[fill=omDef!80, draw=white] table[x=pos, y=mean] {figdata/bachelor-3/20-cost-benefit-analysis/out/overruns_europe.csv};
  \addplot[fill=omThm!70, draw=white] table[x=pos, y=mean] {figdata/bachelor-3/20-cost-benefit-analysis/out/overruns_namerica.csv};
  \addplot[fill=omProp!70, draw=white] table[x=pos, y=mean] {figdata/bachelor-3/20-cost-benefit-analysis/out/overruns_other.csv};
  \addplot[fill=black!35, draw=white] table[x=pos, y=mean] {figdata/bachelor-3/20-cost-benefit-analysis/out/overruns_all.csv};
  \legend{Europe, North America, Other, All regions}
\end{axis}
\end{tikzpicture}
\end{openingchart}

% figdata: overruns_all
Across 258 projects, actual construction costs exceeded the estimates on which the
decision to build was made by 27.6~\% on average: 44.7~\% for rail, 33.8~\% for bridges and
tunnels, 20.4~\% for roads. Costs were underestimated in almost nine projects out of
ten. % ledger: b3cba:fhb.t1
Errors with no systematic cause would scatter around zero; these lean one way,
whether from optimism or from strategic misstatement.

\begin{recall}
The grades 10--12 volume defined present value and discounting. The Year 1 volume
defined the compensating and equivalent variations, and the Year 2 volume the
compensation tests and expected utility. Chapter 5 of this volume estimated the
value of travel time from commuters' choices, and chapter 17 defined the \omterm{def:b3:efficiency-cost-of-taxation:mcpf}{marginal
cost of public funds}.
\end{recall}

\begin{notation}[Symbols of this chapter]\label{not:b3:cost-benefit-analysis:symbols}
$B_t$ and $C_t$ are benefits and costs in year $t$; $\rho$ a discount rate; $p$ a risk of death;
$\mathrm{NPV}$ the net present value.
\end{notation}

\section{Willingness to pay and accounting prices}\label{sec:b3:cost-benefit-analysis:wtp}

\begin{definition}[Cost--benefit analysis; accounting price]\label{def:b3:cost-benefit-analysis:cba}
\index{cost--benefit analysis}\index{accounting price}
\emph{Cost--benefit analysis} values all the effects of a project in money, as the
sum over the people affected of what they would pay for the gains or accept for the
losses, discounts them to the present, and accepts the project if the net present
value is positive. An \emph{accounting price} is the value of a resource used in such an
appraisal when its market price does not measure its cost to society.
\end{definition}

\begin{proposition}[The net-present-value rule and compensation]\label{prop:b3:cost-benefit-analysis:kaldor}
If each person's gain is measured by the compensating variation and money can be
transferred between people at no cost, a positive net present value means the
gainers could compensate the losers and remain better off.
\end{proposition}

\begin{proof}
The compensating variation of each gainer is the most she would pay, and of each
loser the least he would accept, for the project. A positive sum means the gainers'
total willingness to pay exceeds the losers' total compensation: paying each loser
his compensation leaves every gainer at least as well off as without the project,
and some strictly better.
\end{proof}

The rule does not say the compensation is paid; it applies the Year 2 volume's
compensation test, and inherits its limits. Where transfers are costly, a dalen
raised to pay for a project costs society the \omterm{def:b3:efficiency-cost-of-taxation:mcpf}{marginal cost of public funds} of
chapter 17, and the benefits must clear that bar.

\omterm{def:b3:cost-benefit-analysis:cba}{Accounting prices} correct market prices that do not measure costs. A worker hired
from unemployment costs the economy less than her wage: the output she would
otherwise have produced, plus the value of the leisure she gives up. Imported
steel bought under a tariff costs the economy its world price, not the price paid
inside the tariff wall; the tariff is a transfer to the government.

\begin{definition}[Marginal value of public funds]\label{def:b3:cost-benefit-analysis:mvpf}
\index{marginal value of public funds}
The \emph{marginal value of public funds} of a programme is its beneficiaries'
willingness to pay for it divided by its net cost to the government, counting the
taxes its behavioural effects raise or lose.
\end{definition}

The ratio compares programmes of any kind on one scale: a programme whose effects
raise future tax revenue enough to pay for itself has a net cost of zero or less, and
an infinite value.

\begin{record}[The marginal value of public funds of 133 policies]
Hendren and Sprung-Keyser computed the ratio for 133 changes in United States
policy over half a century --- \omterm{def:b3:social-insurance:social-insurance}{social insurance}, education and training, taxes and
cash transfers, transfers in kind --- from the published estimates of each policy's
effects. Direct spending on the health and education of low-income children had the
highest values, above 5 on average. Policies aimed at adults were
generally between 0.5 and 2. % ledger: b3cba:hsk.mvpf
Each value rests on the studies it is built from, and inherits their grades.
\end{record}

\section{Discounting}\label{sec:b3:cost-benefit-analysis:discounting}

\begin{definition}[Social discount rate]\label{def:b3:cost-benefit-analysis:discount-rate}
\index{social discount rate}
The \emph{social discount rate} is the rate used to convert future benefits and costs
into present values in a public appraisal: $\mathrm{NPV} = \sum_t (B_t - C_t)/(1 + \rho)^t$.
\end{definition}

\begin{omchart}{Present value of 1 dalen received in year $t$, at discount rates of
1.4~\%, 3~\% and 7~\% a year.}{Our computation (\code{make\_pv.py})}
\begin{tikzpicture}
\begin{axis}[omrecord, width=0.58\linewidth, height=0.36\linewidth, xmin=0, xmax=100,
  ymin=0, ymax=1.05, xlabel={Year}, ylabel={Present value}, legend pos=north east,
  legend style={font=\footnotesize}]
  \addplot[omDef, thick, mark=none] table[x=t, y=r14] {figdata/bachelor-3/20-cost-benefit-analysis/out/pv.csv};
  \addplot[black, thick, mark=none] table[x=t, y=r3] {figdata/bachelor-3/20-cost-benefit-analysis/out/pv.csv};
  \addplot[omThm, thick, dashed, mark=none] table[x=t, y=r7] {figdata/bachelor-3/20-cost-benefit-analysis/out/pv.csv};
  \legend{1.4~\%, 3~\%, 7~\%}
\end{axis}
\end{tikzpicture}
\end{omchart}

% figdata: pv_points
A dalen in 100 years is worth 0.25 today at 1.4~\%, 0.05 at 3~\% and 0.001 at 7~\%. For a
project whose benefits come in a decade, the choice among these rates matters
little; for one whose benefits come in a century --- a climate policy, a nuclear
waste store --- it decides the answer. The rate has two candidate sources: the return
that capital earns at the margin, which is what the funds would yield elsewhere,
and the rate at which society is willing to trade consumption now for consumption
later. They differ by the taxes and risk premiums between them, and choosing
between them is a choice the evidence does not make.

\begin{dated}{2026}
HM Treasury's appraisal guidance for the United Kingdom discounts at 3.00~\% a year
in real terms for years 1 to 30, 2.50~\% for years 31 to 75 and 2.25~\% from year 76
(Green Book, revised 28 October 2026). % ledger: b3cba:hmt.stpr
The United States Environmental Protection Agency values a statistical life at 7.4
million dollars of 2006, updated to the prices of each analysis. % ledger: b3cba:epa.vsl
The chart's 1.4~\%, 3~\% and 7~\% span the rates in use and in debate.
\end{dated}

\begin{proposition}[Declining certainty-equivalent rates]\label{prop:b3:cost-benefit-analysis:weitzman}
If the discount rate is $\rho_1$ or $\rho_2 > \rho_1$ for ever, each with probability one half, the
certainty-equivalent rate $R(t)$, defined by $e^{-R(t)t} = \tfrac12(e^{-\rho_1t} + e^{-\rho_2t})$, falls with the horizon
$t$ and tends to $\rho_1$.
\end{proposition}

\begin{proof}
Write the right side as $\tfrac12e^{-\rho_1t}\big(1 + e^{-(\rho_2 - \rho_1)t}\big)$. Then $R(t) = \rho_1 + \big(\ln 2 - \ln(1 + e^{-(\rho_2 - \rho_1)t})\big)/t$. At $t \to 0$ the
expected discount factor's slope gives $R = (\rho_1 + \rho_2)/2$; as $t$ grows, $e^{-(\rho_2 - \rho_1)t} \to 0$ and the second
term is $\ln 2/t \to 0$. In between, $R(t)$ falls because the state with the low rate carries a
growing share of the expected discount factor.
\end{proof}

Uncertainty about the rate thus argues for lower rates at long horizons, because
it is the expected discount factor, not the expected rate, that weighs a future
dalen.

\section{The value of a statistical life}\label{sec:b3:cost-benefit-analysis:vsl}

\begin{definition}[Value of a statistical life]\label{def:b3:cost-benefit-analysis:vsl}
\index{value of a statistical life}
The \emph{value of a statistical life} is the willingness to pay for a small reduction
in the risk of death divided by that reduction: what a group would pay, in total,
for measures that prevent one death among them, none of whom knows who it will be.
\end{definition}

\begin{proposition}[The value of a statistical life from expected utility]\label{prop:b3:cost-benefit-analysis:vsl}
A person with wealth $w$, risk of death $p$ and utility $u(w)$ if alive (zero if dead)
will pay $dw$ for a reduction $dp$ in the risk with $dw/dp = u(w)/\big((1 - p)u'(w)\big)$: the \omterm{def:b3:cost-benefit-analysis:vsl}{value of a
statistical life}.
\end{proposition}

\begin{proof}
Expected utility is $(1 - p)u(w)$. Holding it constant, $-u(w)\,dp + (1 - p)u'(w)\,dw = 0$, so a fall in $p$ of
$dp$ is worth $dw = u(w)\,dp/\big((1 - p)u'(w)\big)$.
\end{proof}

The value rises with wealth (through $u(w)/u'(w)$), so richer people pay more for the
same safety: one reason agencies use one value for everyone rather than one per
income.

\begin{model}{Hedonic wages and the value of a statistical life}
\assumes Workers know the fatal risks of jobs and choose among them in competitive
labour markets; employers pay more for riskier jobs to attract workers.
\predicts Wages rise with fatal risk, other job attributes held constant; the
premium per unit of risk is the marginal worker's \omterm{def:b3:cost-benefit-analysis:vsl}{value of a statistical life}.
\tested A review of more than 60 studies of mortality risk premiums in ten
countries finds the premium, and finds that the \omterm{def:b3:cost-benefit-analysis:vsl}{value of a statistical life} rises
with income with an elasticity of about 0.5 to 0.6. % ledger: b3cba:va
\breaks Workers misperceive small risks; people who mind risk least sort into
dangerous jobs, so the premium measures their value, not the population's; and
risky jobs differ from safe ones in attributes no data record, which bias the
premium either way.
\end{model}

\begin{measure}[The value of a statistical life in practice]
A wage premium of 1{,}000 dalens a year for an extra fatal risk of 1 in 10{,}000 a year
implies a value of 10 million dalens. Agencies publish one value for their
appraisals and update it with prices and incomes, as the dated box above shows for
one of them; the value of a project's safety benefits is the deaths it avoids times
that value.
\end{measure}

\section{Risk and revealed or stated preferences}\label{sec:b3:cost-benefit-analysis:preferences}

\begin{definition}[Hedonic wage method; travel-cost method; contingent valuation]\label{def:b3:cost-benefit-analysis:valuation}
\index{hedonic wage method}\index{travel-cost method}\index{contingent valuation}
The \emph{hedonic wage method} infers the value of a job attribute, such as safety,
from the wage differences it commands. The \emph{travel-cost method} infers the value
of a site from what visitors spend in money and time to reach it. \emph{Contingent
valuation} asks people in a survey what they would pay for a good that is not sold.
\end{definition}

Risky outcomes are valued at their expected values when the risks are spread
over many people or many projects, as most public risks are; a project whose
failure would hit the whole economy at once carries a cost of risk on top.

\begin{debate}[Can a survey value what is not sold?]
Some goods --- a wilderness nobody visits, a species --- have no market and leave no
trace in behaviour, so only a survey can value them.
\begin{itemize}
\item One survey of the evidence finds three standing problems: answers to
hypothetical questions overstate what people pay when the payment is real; what
people would pay to gain a good differs widely from what they would accept to lose
it; and values do not grow with the scope of the good (saving one lake is valued
like saving ten). % ledger: b3cba:cv.debate
\emph{Grade 2}: experiments comparing hypothetical with real payments, but for
goods unlike those the method is used for.
\item Another holds that well-designed surveys give usable values where nothing
else does, as for the damage of the 1989 Exxon Valdez oil spill, and that leaving
such goods out of an appraisal values them at zero, a worse error.
\emph{Grade 4}: an argument about the alternative, without a test of the values.
\end{itemize}
\end{debate}

\section{Forecasts against outcomes}\label{sec:b3:cost-benefit-analysis:forecasts}

\begin{definition}[Optimism bias; reference-class forecasting]\label{def:b3:cost-benefit-analysis:optimism-bias}
\index{optimism bias}\index{reference-class forecasting}
\emph{Optimism bias} is the systematic tendency of project forecasts to understate
costs and overstate benefits. \emph{Reference-class forecasting} corrects a forecast
by the distribution of outcomes against forecasts in a class of past projects of
the same kind.
\end{definition}

% figdata: overruns_europe, overruns_namerica, overruns_other
The overruns of the opening chart differ by region: 34.2~\% for rail in Europe, 40.8~\%
in North America and 64.6~\% elsewhere; for roads, 22.4~\% in Europe and 8.4~\% in North
America. % ledger: b3cba:fhb.t2
A reference-class correction raises a rail estimate by the class's average overrun
--- or, to be safe with a given probability, by a high percentile of the class's
distribution --- before the appraisal. It does not explain why the forecasts were
low; it removes the bias whatever its cause.

\begin{record}[Forecasts of demand]
Benefits are forecast too, and the same authors audited the traffic forecasts of
210 transport projects in 14 countries. In nine rail projects out of ten, passenger
numbers had been overestimated, by 106~\% on average: the forecast was about twice
what came. For half of the road projects, traffic differed from its forecast by
more than 20~\% either way, with no lean. Neither kind of forecast had become more
accurate over thirty years. % ledger: b3cba:fsb.demand
Rail projects thus combined costs a half above their estimates with half the
forecast passengers: if benefits scale with passengers, an appraisal built on
both forecasts overstated the benefit--cost ratio about threefold
($1.447 \times 2.06 = 2.98$).
\end{record}

\section{Exercises}\label{sec:b3:cost-benefit-analysis:exercises}

\begin{exercise}[$\star$]\label{exo:b3:cost-benefit-analysis:1}
In the opening chart, which type of project overruns most, and where? What share
of projects overran?
\end{exercise}

\begin{exercise}[$\star$]\label{exo:b3:cost-benefit-analysis:2}
A bridge costs 100 million dalens today and yields 6 million a year for 30 years,
from next year. Find its net present value at 3~\% and 7~\%.
\end{exercise}

\begin{exercise}[$\star$]\label{exo:b3:cost-benefit-analysis:3}
Workers accept a fatal risk higher by 1 in 10{,}000 a year for 1{,}000 dalens more pay a
year. Find the \omterm{def:b3:cost-benefit-analysis:vsl}{value of a statistical life}. What is it if the premium is 800?
\end{exercise}

\begin{exercise}[$\star$]\label{exo:b3:cost-benefit-analysis:4}
What data on forecasts would distinguish random errors from \omterm{def:b3:cost-benefit-analysis:optimism-bias}{optimism bias}?
\end{exercise}

\begin{exercise}[$\star\star$]\label{exo:b3:cost-benefit-analysis:5}
A project hires 100 workers at 30{,}000 dalens a year. Sixty would otherwise be
unemployed, valuing their leisure at 8{,}000 a year; forty would leave other jobs where
they produced 30{,}000. Find the accounting cost of the labour.
\end{exercise}

\begin{exercise}[$\star\star$]\label{exo:b3:cost-benefit-analysis:6}
The discount rate is 1~\% or 5~\%, with equal probability, for ever. Compute the
certainty-equivalent rate at 10, 50 and 200 years.
\end{exercise}

\begin{exercise}[$\star\star$]\label{exo:b3:cost-benefit-analysis:7}
A climate-type project costs 10 today and yields 100 in year 100. For which discount
rates is it worth doing?
\end{exercise}

\begin{exercise}[$\star\star$]\label{exo:b3:cost-benefit-analysis:8}
A rail project is forecast to cost 2 billion dalens. Correct the forecast with the
rail averages of the opening chart for Europe and for all regions, and say what a
correction to the 80th percentile would need.
\end{exercise}

\begin{exercise}[$\star\star$]\label{exo:b3:cost-benefit-analysis:9}
What evidence would show that stated willingness to pay is biased?
\end{exercise}

\begin{exercise}[$\star\star\star$]\label{exo:b3:cost-benefit-analysis:10}
Prove the declining-rate proposition, and show that $R(t)$ is strictly decreasing.
\end{exercise}

\begin{exercise}[$\star\star\star$]\label{exo:b3:cost-benefit-analysis:11}
Simulate 5{,}000 workers whose log wage depends on schooling, experience and a job's
fatal risk, with a premium of 0.02 log points per extra death in 10{,}000, at a mean
wage of 30{,}000 dalens. Estimate the hedonic wage regression and the \omterm{def:b3:cost-benefit-analysis:vsl}{value of a
statistical life}, then add an unrecorded job attribute correlated with risk and
show the bias.
\end{exercise}

\begin{exercise}[$\star\star\star$]\label{exo:b3:cost-benefit-analysis:12}
Programme A costs 100, raises 20 in later taxes, and its beneficiaries would pay 120
for it. Programme B costs 100, loses 10 in taxes through lower work, and its
beneficiaries would pay 90. Find each one's \omterm{def:b3:cost-benefit-analysis:mvpf}{marginal value of public funds} and rank
them.
\end{exercise}

\section{Problem: A New Bridge for Brenmoor}\label{sec:b3:cost-benefit-analysis:problem}

\begin{problem}[{Weekend problem --- A new bridge for Brenmoor}]\label{pb:b3:cost-benefit-analysis:bridge}
Brenmoor, Brenland's capital, plans a bridge forecast to cost 120, 160 and 120
million dalens in years 0, 1 and 2. It opens in year 3 and is appraised over 40 years
of use, with maintenance of 4 million a year. It will carry 30{,}000 vehicle trips a
day, each saving 10 minutes for 1.2 occupants, valued at 13.1 dalens an hour (the
value of travel time of chapter 5). It is expected to avoid one road death a year,
valued at 3 million dalens. % figdata: bridge_stats, bridge_npv

\medskip\noindent\textbf{Part I --- the benefits.}
\begin{enumerate}
\item Compute the hours saved in a year.
\item Value them.
\item Why use commuters' own value of time rather than their wage?
\item Compute the safety benefit.
\item Compute the net annual benefit after maintenance.
\item Which effects of a bridge does this list leave out?
\end{enumerate}

\medskip\noindent\textbf{Part II --- safety.}
\begin{enumerate}[resume]
\item How would a hedonic-wage study produce the 3 million?
\item What would the safety benefit be with a value of 6 million?
\item Why should every death avoided carry the same value, whoever travels?
\item Workers in risky jobs may mind risk least. Which way does that bias the value?
\end{enumerate}

\medskip\noindent\textbf{Part III --- discounting.}
\begin{enumerate}[resume]
\item Write the net present value as a sum over years.
\item Compute it at 2~\%, 3.5~\% and 7~\%.
\item Find the discount rate at which it is zero.
\item Explain why the NPV falls with the rate for this project.
\item Which rate would you use, and why does the choice matter here?
\item The net-present-value chart below plots both cases of Part IV. Read off where
each crosses zero.
\end{enumerate}

\begin{center}
\begin{tikzpicture}
\begin{axis}[omrecord, width=0.55\linewidth, height=0.34\linewidth, xmin=0, xmax=10,
  ymin=-400, ymax=800, xlabel={Discount rate (\%)},
  ylabel={NPV (million dalens)}, legend pos=north east,
  legend style={font=\footnotesize}]
  \addplot[omDef, thick, mark=none] table[x=rate, y=npv] {figdata/bachelor-3/20-cost-benefit-analysis/out/bridge_npv.csv};
  \addplot[omThm, thick, dashed, mark=none] table[x=rate, y=npv_c] {figdata/bachelor-3/20-cost-benefit-analysis/out/bridge_npv.csv};
  \draw[guide] (axis cs:0,0) -- (axis cs:10,0);
  \legend{forecast cost, corrected cost}
\end{axis}
\end{tikzpicture}
\end{center}

\medskip\noindent\textbf{Part IV --- a reference-class correction.} Bridges and tunnels
overran their estimates by 33.8~\% on average.
\begin{enumerate}[resume]
\item Correct the construction cost.
\item Recompute the net present value at 2~\%, 3.5~\% and 7~\%.
\item Find the new rate at which it is zero.
\item Should traffic forecasts be corrected too? What would that do?
\item The money is raised by taxes with a \omterm{def:b3:efficiency-cost-of-taxation:mcpf}{marginal cost of public funds} of 1.3. How
does that change the test?
\item Why might the forecasters' estimate be low in the first place?
\item What would make the correction itself wrong for this bridge?
\item Who gains from the bridge, and who pays?
\item State the result: the bridge's net present value after the correction, at
3.5~\%, and the discount rate at which it turns negative.
\end{enumerate}
\end{problem}
